\documentclass[portrait,a0]{a0poster}
\usepackage{ascmac}
\usepackage[dvips]{graphicx,epsfig,color}
\usepackage{multicol}
\usepackage{multirow}
\setlength{\columnsep}{20pt}
\pagestyle{empty}
\def\spaceV{\ \ \ \ \ }
\def\spaceX{\ \ \ \ \ \ \ \ \ \ }
\def\spaceXV{\ \ \ \ \ \ \ \ \ \ \ \ \ \ \ }
\def\spaceXX{\ \ \ \ \ \ \ \ \ \ \ \ \ \ \ \ \ \ \ \ }
\def\ColoraTProCy{\color[cmyk]{0.7,0.7,0,0}}
\def\ColorbTAnaeroI{\color[cmyk]{1,0,1,0}}
\def\ColorcTAnaeroII{\color[cmyk]{0.5,0,0.5,0}}
\def\ColordTAnaeroIII{\color[cmyk]{0.2,0,0.2,0}}
\def\ColoreTAnfMethI{\color[cmyk]{1,0,0,0}}
\def\ColorfTAnfMethII{\color[cmyk]{0.4,0,0,0}}
\def\ColorgTPNifI{\color[cmyk]{0,1,1,0}}
\def\ColorhTPNifII{\color[cmyk]{0,0.7,0.7,0}}
\def\ColoriTPNifIII{\color[cmyk]{0,0.75,0,0}}
\def\ColorjTPNifIV{\color[cmyk]{0.5,1,0.5,0}}
\def\ColorkBK{\color[cmyk]{0,0,0,1}}

\begin{document}
\linespread{1.4}\selectfont
%%----- Title -----%%
\begin{center}
\color[cmyk]{1,1,0,0}
\fontsize{128}{20}\selectfont
Self-Organizing Clustering \\ $BF|K\8l(BOK
\end{center}
\begin{center}
\color[cmyk]{1,1,0,0}
\fontsize{56}{20}\selectfont
-- A Novel Non-Hierarchical Method for Clustering Large Amount of DNA Sequences --
\end{center}
%%----- /Title -----%%
\vspace{9mm}
%%----- Author(s) and Affiliation(s) -----%%
\color[cmyk]{1,1,0,0.4}
\fontsize{30}{20}\selectfont
\begin{center} \begin{tabular}[t]{ccccc}
\bf Kou Amano\footnotemark[1] \footnotemark[2] & \hspace{20mm} & \bf Hidemitsu Nakamura\footnotemark[2] & \hspace{20mm} & \bf Hiroaki Ichikawa\footnotemark[2] \\
\tt amano@ulis.ac.jp & \  & \tt hide7087@affrc.go.jp & \  & \tt hichkw@affrc.go.jp \\
\end{tabular} \end{center}
\begin{center} \begin{tabular}{rl}
\tt \footnotemark[1] & \tt University of Library and Information Science, 1-2 Kasuga, Tsukuba, Ibaraki 305-8550 Japan \\
\tt \footnotemark[2] & \tt National Institute of Agrobiological Sciences, 2-1-2 Kannondai, Tsukuba, Ibaraki 305-8602 Japan \\
\end{tabular} \end{center}
%%----- /Author(s) and Affiliation(s) -----%%
\vspace{10mm}
%%----- Multi Column Region -----%%
\linespread{1.6}\fontsize{32}{20}\selectfont
\begin{multicols}{3}
\color[cmyk]{0,0,0,1}
\section{Introduction}
\

\vspace{-40pt} To cluster and characterize DNA sequences with focusing on the oligonucleotide frequency, we developed a novel method and program package designated as Self-Organizing Clustering (SOC).

Being based on Self-Organizing Map (SOM)\cite{Ikemura-2001}, the algorithm of SOC made use of K-means with modification.

\vspace{-5mm}\section{Aim of SOC}
\

\vspace{-40pt} The SOC is a non-hierarchical clustering method to cluster the samples with multi-dimensional data into several groups.

Among several non-hierarchical methods for clustering, the aim of SOC is to cluster a huge amount of specimens each comprised of a long stretch of nucleotide sequence.

\vspace{-5mm}\section{Existing tools vs. SOC}

%%----- Table 1 -----%%
\begin{center}
\linespread{0.9}\selectfont 
\begin{tabular}{cccp{80mm}}
\multicolumn{4}{l}{\textit{Table 1. Comparison of clustering tools}}\\
\hline
\textgt{program} & \textgt{cost} & \textgt{performance} & \hspace{10mm}\textgt{features}\\
\hline
\\
\color[cmyk]{0,1,1,0}\textbf{SOC} & \color[cmyk]{0,1,1,0}Free & \color[cmyk]{0,1,1,0}Fast & \small \begin{itemize}
	\vspace{-9mm}
	\color[cmyk]{0,1,1,0}\item[] Specialized for clustering.
	\color[cmyk]{0,1,1,0}\item[] Not using neural-networks.
	\color[cmyk]{0,1,1,0}\item[] Not based on alignment. \end{itemize}\\
SOM & Free  & Fast & \small \begin{itemize}
	\vspace{-9mm}
	\item[] Original purpose is simulation of learning and memory. \end{itemize}\\
BL-SOM & Expensive & Fast & \small \begin{itemize}
	\vspace{-9mm}
	\item[] Original purpose is data mining.
	\item[] Makes two-dimensional map.
	\item[] Clusters are separated by borders on map. \end{itemize}\\
CLUSTALX/W & Free & Slow & \small \begin{itemize}
	\vspace{-9mm}
	\item[] Based on alignment. \end{itemize}\\
\hline
\\
%\multicolumn{4}{l}{\textit{Table 1. Comparison of clustering tools}}\\[-6mm]
\end{tabular}
\end{center}
%%----- /Table 1 -----%%

\vspace{-3mm}\section{How does the SOC Work?}
\

\vspace{-40pt} SOC works as follows.
\vspace{3mm}
\begin{flushleft}
\linespread{1.8}\selectfont
\begin{minipage}{245mm}
\begin{center} \end{center}
\begin{tabular}{l}\begin{screen}\linespread{1.2}\selectfont
1. Read data of  sample-nodes.
\end{screen}\end{tabular}
\begin{center}\vspace{-5mm}$|$\vspace{-3mm}\end{center}
\begin{tabular}{l}\begin{screen}\linespread{1.2}\selectfont
2. \hspace{-2mm}Place cluster-nodes at random \hspace{1mm}on \hspace{1mm}the same space as sample-nodes.
\end{screen}\end{tabular}
\begin{center}\vspace{-5mm}$|$\vspace{-3mm}\end{center}
\begin{tabular}{l}\begin{screen}\linespread{1.2}\selectfont
3. Assign each sample-node to the nearest cluster-node.
\end{screen}\end{tabular}
\put(-5,5){\line(1,0){10}\line(0,-1){428}}\vspace{-402pt}
\begin{center}\vspace{-5mm}$|$\vspace{-3mm}\end{center}
\begin{tabular}{l}\begin{screen}\linespread{1.2}\selectfont
4. \hspace{-5mm}Calculate each coordinate for a centroid of a set of sample-nodes which are assigned to the particular cluster-node.
\end{screen}\end{tabular}
\begin{center}\vspace{-5mm}$|$\vspace{-3mm}\end{center}
\begin{tabular}{l}\begin{screen}\linespread{1.2}\selectfont
5. Move each cluster-node toward the given centroid.
\end{screen}\end{tabular}
\begin{center}\vspace{-5mm}$|$\vspace{-3mm}\end{center}
\begin{tabular}{l}\begin{screen}\linespread{1.2}\selectfont
6. \hspace{-6mm}Calculate \hspace{2mm}the \hspace{2mm}distance between \hspace{1mm}given \hspace{1mm}two cluster-nodes.
\end{screen}\end{tabular}
\begin{center}\vspace{-5mm}$|$\vspace{-3mm}\end{center}
\begin{tabular}{l}\begin{screen}\linespread{1.2}\selectfont
7. Unify a pair of cluster-nodes with the shortest distance, Unless the distance is longer than a threshold value given by the user.
\end{screen}\end{tabular}
\put(-5,5){\line(1,0){10}}
\begin{center}\vspace{-3mm}\end{center}
\vspace{-5mm}
\begin{tabular}{l}\linespread{1.2}\selectfont
\ Repeat steps 3 to 7.
\end{tabular}
\end{minipage}

\end{flushleft}

\vspace{-5mm}\section{Clustering genes from 17 bacteria}
\

\vspace{-40pt} We used a total of 800 genes extracted from 17 bacterial genomes as a set of test samples.
The lengths of the DNA sequences were 2.0 to 2.2 kb. 
When trinucleotide frequency was analyzed, the samples were divided into eight clusters (Table 2).
Genes derived from a given bacterial species were classified into a particular cluster with high accuracy.

\section{Clustering \textit{nifH} genes}
\

\vspace{-40pt} Nitrogen fixation by the microorganisms living in termite's gut is one of the crucial aspects of symbiosis, since termites usually thrive on a nitrogen-poor diet.
 
The \textit{nifH} gene encodes dinitrogenase reductase and has often been used as a molecular marker for nitrogen-fixing microorganisms. Many different \textit{nifH} genes, which might be originated from several microbial species, are isolated from the gut of a single termite. 
Ohkuma {\em et al.}\cite{Ohkuma} isolated \textit{nifH} genes from DNA preparations extracted from guts of nine termite species, and clustered them into 10 major groups (Table 3). They used CLUSTALW and PHYLIP program package to draw the phylogenetic trees. According to the classification by Ohkuma {\em et al.}, gene shown by the same color belong to the same cluster.

We clustered 145 \textit{nifH} genes by the SOC. Interestingly, many genes in the same group classified by the CLUSTALW were dispersed into two or more groups.

\section{Conclusions}

\begin{itemize}
\item Though some clusters were consisted of genes from two or more species, the SOC grouped DNA sequences from a given bacterial species into a particular cluster with high accuracy in the classification of 800 genes from 17 bacterial genomes.
In progress is the further clustering of such clusters by using the SOC hierarchically.
\item We classified \textit{nifH} genes from the gut of diverse termites by the SOC. The grouping pattern of the SOC was quite different from that of the CLUSTALW. 
Although further analysis would be needed to explain what the grouping pattern of the SOC means, 
our results strongly indicate that the SOC might be of use to classify a complex pool of DNA sequences by a novel strategy.
\end{itemize}
%%----- Table 2 -----%%
\begin{center}
\color[cmyk]{0,0,0,1}
\begin{tabular}{|c|c|c|}
\multicolumn{3}{l}{\textit{Table 2. Clustering 800 genes from 17 bacteria.}}\\
\hline
cluster ID&species&allocation(\%)\\
\hline
I&\color[cmyk]{0,0.7,0.7,0}\textit{\textbf{Mesorhizobium loti}}&\color[cmyk]{0,0.7,0.7,0}98\\
\hline
II&\color[cmyk]{0,0,0,1}\textit{\textbf{Aeropyrum pernix}}&\color[cmyk]{0,0,0,1}4\\
&\begin{tabular}{c}
	\color[cmyk]{0,0.7,0.7,0}\textit{\textbf{Agrobacterium tumefaciens}} \\[-3mm] \color[cmyk]{0,0.7,0.7,0}\textbf{C58}\end{tabular}&\color[cmyk]{0,0.7,0.7,0}97\\
&\color[cmyk]{0,0.7,0.7,0}\textit{\textbf{Bifidobacterium longum}}&\color[cmyk]{0,0.7,0.7,0}97\\
&\color[cmyk]{0,0.7,0.7,0}\textit{\textbf{Brucella melitensis}}&\color[cmyk]{0,0.7,0.7,0}96\\
&\color[cmyk]{0,0.7,0.7,0}\textit{\textbf{Buchnera aphidicola}}&\color[cmyk]{0,0.7,0.7,0}94\\
\hline
III&\color[cmyk]{0,0.7,0.7,0}\textit{\textbf{A. pernix}}&\color[cmyk]{0,0.7,0.7,0}96\\
&\color[cmyk]{0,0,0,1}\textit{\textbf{A. tumefaciens}} \textbf{C58}&\color[cmyk]{0,0,0,1}3\\
&\color[cmyk]{0,0,0,1}\textit{\textbf{Bacillus halodurans}}&\color[cmyk]{0,0,0,1}10\\
&\color[cmyk]{0,0,0,1}\textit{\textbf{B. longum}}&\color[cmyk]{0,0,0,1}3\\
&\color[cmyk]{0,0,0,1}\textit{\textbf{B. melitensis}}&\color[cmyk]{0,0,0,1}4\\
&\color[cmyk]{0,0,0,1}\textit{\textbf{B. aphidicola}}&\color[cmyk]{0,0,0,1}6\\
&\begin{tabular}{c}
	\color[cmyk]{0,0,0,1}\textit{\textbf{Chlamydophila pneumoniae}} \\[-3mm] \textbf{AR39}\end{tabular}&\color[cmyk]{0,0,0,1}48\\
&\color[cmyk]{0,0.7,0.7,0}\textit{\textbf{Haemophilus influenzae}}&\color[cmyk]{0,0.7,0.7,0}83\\
&\color[cmyk]{0,1,1,0}\textit{\textbf{Lactobacillus plantarum}}&\color[cmyk]{0,1,1,0}100\\
&\color[cmyk]{0,0,0,1}\textit{\textbf{Leptospira interrogans}}&\color[cmyk]{0,0,0,1}1\\
&\color[cmyk]{0,0,0,1}\textit{\textbf{M. loti}}&\color[cmyk]{0,0,0,1}2\\
\hline
IV&\color[cmyk]{0,0.7,0.7,0}\textit{\textbf{Deinococcus radiodurans}}&\color[cmyk]{0,0.7,0.7,0}98\\
\hline
V&\color[cmyk]{0,0,0,1}\textit{\textbf{C. pneumoniae}} \textbf{AR39}&\color[cmyk]{0,0,0,1}4\\
&\color[cmyk]{0,1,1,0}\textit{\textbf{Clostridium acetobutylicum}}&\color[cmyk]{0,1,1,0}100\\
&\color[cmyk]{0,1,1,0}\textit{\textbf{Fusobacterium nucleatum}}&\color[cmyk]{0,1,1,0}100\\
\hline
VI&\color[cmyk]{0,0.7,0.7,0}\textit{\textbf{B. halodurans}}&\color[cmyk]{0,0.7,0.7,0}90\\
&\color[cmyk]{0,1,1,0}\textit{\textbf{Campylobacter jejuni}}&\color[cmyk]{0,1,1,0}100\\
&\color[cmyk]{0,1,1,0}\textit{\textbf{Chlamydia muridarum}}&\color[cmyk]{0,1,1,0}100\\
&\color[cmyk]{0,0,0,1}\textit{\textbf{C. pneumoniae}} \textbf{AR39}&\color[cmyk]{0,0,0,1}48\\
\hline
VII&\color[cmyk]{0,0,0,1}\textit{\textbf{H. influenzae}}&\color[cmyk]{0,0,0,1}4\\
&\color[cmyk]{0,0.7,0.7,0}\textit{\textbf{L. interrogans}}&\color[cmyk]{0,0.7,0.7,0}99\\
\hline
VIII&\color[cmyk]{0,0,0,1}\textit{\textbf{D. radiodurans}}&\color[cmyk]{0,0,0,1}2\\
&\color[cmyk]{0,1,1,0}\textit{\textbf{Escherichia coli}} \textbf{K12}&\color[cmyk]{0,1,1,0}100\\
&\color[cmyk]{0,0,0,1}\textit{\textbf{H. influenzae}}&\color[cmyk]{0,0,0,1}13\\
\hline
\multicolumn{3}{p{240mm}}{\linespread{0.9}\fontsize{22}{20}\selectfont \color[cmyk]{0,0,0,1}Species shown by \color[cmyk]{0,0.7,0.7,0}red \color[cmyk]{0,0,0,1}color carry more than 80\% of gene allocations on a particular cluster.}\\[-8mm]
%\multicolumn{3}{p{240mm}}{\linespread{0.9}\fontsize{20}{20}\selectfont \color[cmyk]{0,0,0,1}Greek number \color[cmyk]{0,0,0,1}: Cluster by SOC.}\\[-8mm]
%\multicolumn{3}{p{240mm}}{\linespread{0.9}\fontsize{20}{20}\selectfont \color[cmyk]{0,1,1,0}Color \color[cmyk]{0,0,0,1}: Allocated to a cluster with high percentage ($\% \geq 80$). Same color correspond to same bactarial species.}\\[-8mm]
%\multicolumn{3}{p{240mm}}{\linespread{0.9}\fontsize{20}{20}\selectfont \color[cmyk]{0,0,0,1}Black \color[cmyk]{0,0,0,1}: Dispersed to several clusters.}\\
\end{tabular}
\end{center}

%%----- /Table 2 -----%%
\vspace{5mm}
%%----- Table 3 -----%%
\begin{center}
\begin{tabular}{ccp{150mm}}
\multicolumn{3}{p{220mm}}{\textit{Table 3. Clustering 145 genes extracted from the mixed populations in the termite gut.}}\\
\hline
cluster ID&\ &\spaceXX  gene \\
\hline
I&&\fontsize{24}{20}\selectfont 
\ColorbTAnaeroI NTN24,
\ColoriTPNifIII GFN15,
\ColorkBK NTN30
\\
\hline
II&&\fontsize{24}{20}\selectfont 
\ColorgTPNifI GFN12,
\ColorgTPNifI NKN24,
\ColorgTPNifI NTN39,
\ColorgTPNifI OFN16,
\ColorgTPNifI OFN17,
\ColorgTPNifI OFN39,
\ColoriTPNifIII PNN24,
\ColorjTPNifIV NTN19,
\ColorjTPNifIV NTN29,
\ColorjTPNifIV NTN4,
\ColorjTPNifIV PNN26,
\ColorjTPNifIV PNN49,
\ColorkBK CDN10,
\ColorkBK CDN12,
\ColorkBK CDN14,
\ColorkBK NTN7,
\ColorkBK TKY23,
\ColorkBK TKY24,
\ColorkBK TKY26
\\
\hline
III&&\fontsize{24}{20}\selectfont 
\ColorbTAnaeroI TKY15,
\ColorbTAnaeroI TKY7,
\ColorcTAnaeroII CDN1,
\ColorcTAnaeroII CDN3,
\ColorcTAnaeroII CDN7,
\ColorfTAnfMethII GFN18,
\ColorfTAnfMethII GFN7,
\ColorgTPNifI CDN25,
\ColorgTPNifI CDN33,
\ColorgTPNifI CFN9,
\ColorgTPNifI GFN3,
\ColorgTPNifI GFN9,
\ColorgTPNifI TKY1,
\ColorgTPNifI TKY22,
\ColoriTPNifIII GFN16,
\ColorjTPNifIV CDN31,
\ColorjTPNifIV CDN34,
\ColorjTPNifIV NKN12,
\ColorjTPNifIV OFN12,
\ColorjTPNifIV PNN23,
\ColorjTPNifIV TKY19,
\ColorkBK HSN14,
\ColorkBK NKN23
\\
\hline
IV&&\fontsize{24}{20}\selectfont 
\ColoraTProCy TKY17,
\ColorcTAnaeroII GFN25,
\ColorcTAnaeroII HSN10,
\ColorcTAnaeroII HSN11,
\ColorcTAnaeroII HSN16,
\ColorcTAnaeroII HSN20,
\ColorcTAnaeroII NTN2,
\ColorcTAnaeroII NTN21,
\ColorcTAnaeroII TKY9,
\ColorfTAnfMethII HSN19,
\ColorfTAnfMethII HSN9,
\ColoriTPNifIII NKN17,
\ColorjTPNifIV NTN22,
\ColorjTPNifIV NTN32,
\ColorkBK HSN17,
\ColorkBK HSN18,
\ColorkBK HSN2,
\ColorkBK HSN24,
\ColorkBK HSN26,
\ColorkBK HSN4,
\ColorkBK HSN5,
\ColorkBK HSN6,
\ColorkBK NTN38,
\ColorkBK TKY11,
\ColorkBK TKY14
\\
\hline
V&&\fontsize{24}{20}\selectfont 
\ColorcTAnaeroII TKY2,
\ColordTAnaeroIII CDN8,
\ColordTAnaeroIII GFN14,
\ColordTAnaeroIII GFN8,
\ColordTAnaeroIII NKN6,
\ColordTAnaeroIII PNN16,
\ColoreTAnfMethI CDN22,
\ColoreTAnfMethI CDN23,
\ColoreTAnfMethI CDN4,
\ColoreTAnfMethI CDN6,
\ColoreTAnfMethI NKN3,
\ColoreTAnfMethI NKN4,
\ColoreTAnfMethI NKN9,
\ColoreTAnfMethI OFN35,
\ColorfTAnfMethII CDN5,
\ColorgTPNifI GFN19,
\ColorgTPNifI NTN16,
\ColorjTPNifIV PNN48,
\ColorjTPNifIV PNN51,
\ColorkBK CDN24,
\ColorkBK CDN27,
\ColorkBK NKN10,
\ColorkBK NKN14,
\ColorkBK NKN15,
\ColorkBK NKN16,
\ColorkBK NKN21,
\ColorkBK NKN29,
\ColorkBK NKN30,
\ColorkBK NKN8,
\ColorkBK NTN34,
\ColorkBK PNN46
\\
\hline
VI&&\fontsize{24}{20}\selectfont 
\ColorbTAnaeroI CFN21,
\ColorbTAnaeroI CFN3,
\ColorbTAnaeroI CFN4,
\ColorbTAnaeroI CFN6,
\ColorbTAnaeroI TKY16,
\ColorbTAnaeroI TKY6,
\ColorcTAnaeroII NTN6,
\ColorhTPNifII OFN1,
\ColorhTPNifII OFN3,
\ColorhTPNifII PNN31,
\ColorkBK CFN25,
\ColorkBK CFN5,
\ColorkBK CFN7,
\ColorkBK GFN1,
\ColorkBK GFN24,
\ColorkBK TKY13,
\ColorkBK TKY18,
\ColorkBK TKY20,
\ColorkBK TKY4,
\ColorkBK TKY8
\\
\hline
VII&&\fontsize{24}{20}\selectfont 
\ColorcTAnaeroII CFN10,
\ColorcTAnaeroII CFN12,
\ColorcTAnaeroII NTN18,
\ColordTAnaeroIII CFN13,
\ColordTAnaeroIII CFN24,
\ColorhTPNifII PNN18,
\ColorjTPNifIV PNN17,
\ColorjTPNifIV PNN27,
\ColorkBK CFN17,
\ColorkBK NKN19,          
\ColorkBK PNN1,
\ColorkBK PNN47
\\
\hline
VIII&&\fontsize{24}{20}\selectfont 
\ColorbTAnaeroI CFN22,
\ColorbTAnaeroI GFN10,
\ColorbTAnaeroI NKN20, 
\ColorbTAnaeroI NKN5,
\ColorbTAnaeroI TKY3,
\ColorcTAnaeroII NTN36,
\ColoriTPNifIII PNN21, 
\ColoriTPNifIII PNN32,
\ColorjTPNifIV NKN11,
\ColorkBK CDN21, 
\ColorkBK CDN28,
\ColorkBK CFN16
\\
\hline
%\multicolumn{3}{p{220mm}}{\textit{Table 3. Clustering 145 genes extracted from the mixed population in the termite gut.}}\\[-6mm]
%\multicolumn{3}{p{220mm}}{\fontsize{20}{20}\selectfont\ColorkBK Greek number \ColorkBK : Cluster by SOC.}\\[-8mm]
\multicolumn{3}{p{220mm}}{\fontsize{22}{20}\selectfont\ColoraTProCy Blue \ColorkBK : Proteo-Cyano group.}\\[-8mm]
\multicolumn{3}{p{220mm}}{\fontsize{22}{20}\selectfont\ColorbTAnaeroI Deep green \ColorkBK : Anaerobe \textit{nif} group cluster I.}\\[-8mm]
\multicolumn{3}{p{220mm}}{\fontsize{22}{20}\selectfont\ColorcTAnaeroII Green \ColorkBK : Anaerobe \textit{nif} group cluster II.}\\[-8mm]
\multicolumn{3}{p{220mm}}{\fontsize{22}{20}\selectfont\ColordTAnaeroIII Light green \ColorkBK : Anaerobe \textit{nif} group cluster III.}\\[-8mm]
\multicolumn{3}{p{220mm}}{\fontsize{22}{20}\selectfont\ColoreTAnfMethI Deep cyan \ColorkBK : \textit{anf}-methano group cluster I.}\\[-8mm]
\multicolumn{3}{p{220mm}}{\fontsize{22}{20}\selectfont\ColorfTAnfMethII cyan \ColorkBK : \textit{anf}-methano group cluster II.}\\[-8mm]
\multicolumn{3}{p{220mm}}{\fontsize{22}{20}\selectfont\ColorgTPNifI Deep red \ColorkBK : Pseudo \textit{nif} group cluster I.}\\[-8mm]
\multicolumn{3}{p{220mm}}{\fontsize{22}{20}\selectfont\ColorhTPNifII Red \ColorkBK : Pseudo \textit{nif} group cluster II.}\\[-8mm]
\multicolumn{3}{p{220mm}}{\fontsize{22}{20}\selectfont\ColoriTPNifIII Magenta \ColorkBK : Pseudo \textit{nif} group cluster III.}\\[-8mm]
\multicolumn{3}{p{220mm}}{\fontsize{22}{20}\selectfont\ColorjTPNifIV Purple \ColorkBK : Pseudo \textit{nif} group cluster IV.}\\[-8mm]
\multicolumn{3}{p{220mm}}{\fontsize{22}{20}\selectfont\ColorkBK Black \ColorkBK : Except the 10 groups.}\\
\end{tabular}
\end{center}

%%----- /Table 3 -----%%

%%----- Acknowledgments -----%%
\section*{Acknowledgments} \linespread{1.4}\fontsize{25}{20}\selectfont
\

\vspace{-40pt} We thank Dr. N. Onodera (School of Library and Information Science, University of Tsukuba),
Dr. T. Ikemura (National Institute of Genetics),
Dr. T. Maeshiro (Research Center for Knowledge Communities, University of Tsukuba),
Mr. T. Uchida (NABE International),
and Mr. T. Kozuki (Xanagen) 
for helpful suggestions.

We also thank Dr. M. Ohkuma (RIKEN) for providing DNA data of \textit{nifH} genes. \normalfont
%%----- /Acknowledgments -----%%

%%----- References -----%%
\vspace{-6mm}
\begin{thebibliography}{10}
%\bibitem{Ikemura-1998}
%Ikemura, T.
%Identification of gene regulatory networks by strategic gene
%disruptions and gene overexpressions,
%{\em Proc. 9th ACM-SIAM Symp. Discrete Algorithms}, 695--702, 1998.

\bibitem{Ikemura-2001}
Kanaya, S., Kinouchi, M., Abe, T., Kudo, Y., Yamada, Y., Nishi, T., Mori, H., and Ikemura, T.,
Analysis of codon usage diversity of bacterial genes with a self-organizing map(SOM): characterization of horizontally transferred genes with emphasis on the {\it E. coli} O157 genome,
{\em Gene}, 276:89--99, 2001.
%
%\bibitem{Koho}
%Kohonen, T.
%{\em Self-organizing maps}, Springer, 1995.

\bibitem{Ohkuma}
Ohkuma, M., Noda, S., and Kudo, T.,
Phylogenetic Diversity of Nitrogen Fixation Genes in the Symbiotic Microbial Community in the Gut of Termites,
{\em Applied and Environmental Microbiology}, 65(11):4926--4934, 1999.
\end{thebibliography}
%%----- /References -----%%
\vspace{10mm}
%%----- Date and JSBI logo-----%%
\begin{flushright}
\color[cmyk]{0,0,0,1}\fontsize{28}{20}\selectfont \textgt{GIW2003 15-17 Dec.}
\epsfig{file=JSBi2-bk.ps,height=8mm}
\end{flushright}
%%----- /Data and JSBi Logo -----%%
\end{multicols}
%%----- /Multi Column Region -----%%

%%----- Title Logo -----%%
%\vspace{20mm}
\begin{center}
\epsfig{file=SOC_TITLE_LOGO.ps,height=30mm}
\end{center}
%%----- /Title Logo -----%%

%%----- NIAS Logo -----%%
%vspace{20mm}
\begin{center}
\epsfig{file=nias_logo.ps,height=45mm}
\end{center}
%%----- /NIAS Logo -----%%

%%----- For print all image -----%%
\begin{center}
\epsfig{file=WhiteSpace.ps,height=20mm}
\end{center}
%%----- For print all image -----%%

\end{document}
